%==============================================================================
%  Musa Ibne Mannan - Resume
%  Built for ATS (applicant tracking system) parsing:
%    - T1 font encoding + glyphtounicode, so PDF text extracts as real
%      characters instead of mangled ligatures
%    - hyphenation disabled, so no word is split across lines with a hyphen
%    - standard section names ("Summary", "Education", "Experience", "Skills",
%      "Publications", "Projects") that parsers recognize
%    - single column, no tables, no text boxes, no header/footer content
%    - month-year date format
%  Compile with pdflatex.
%==============================================================================

\documentclass[11pt, letterpaper]{article}

\usepackage[margin=0.5in]{geometry}
\linespread{0.97}  % 3%% tighter leading, holds the resume to two pages
\usepackage[T1]{fontenc}          % T1 encoding: required for clean text extraction
\usepackage[utf8]{inputenc}
\usepackage{mathptmx}             % Times text and math (replaces the obsolete `times`)
\usepackage{enumitem}
\usepackage{titlesec}
\usepackage{xcolor}
\usepackage[hidelinks]{hyperref}

% --- ATS: embed a ToUnicode map so copy/paste and parsing return real text ----
\input{glyphtounicode}
\pdfgentounicode=1

% --- ATS: never hyphenate; a broken word confuses keyword matching -----------
\hyphenpenalty=10000
\exhyphenpenalty=10000
\tolerance=2000
\emergencystretch=3em

\hypersetup{
  pdftitle={Musa Ibne Mannan - Resume},
  pdfauthor={Musa Ibne Mannan},
  pdfsubject={Mechanical Engineering, Failure Analysis, Finite Element Analysis},
  pdfkeywords={failure analysis, root cause analysis, finite element analysis,
    Abaqus, FEA, materials characterization, mechanical testing, SEM, EDX, XPS,
    XRD, Python, MATLAB, SolidWorks, AutoCAD, CATIA, CVD, Raman spectroscopy,
    forensic engineering, technical consulting, mechanical engineering}
}

\definecolor{accent}{RGB}{0, 51, 102}

\titleformat{\section}{\large\bfseries\color{accent}}{}{0em}{}[\titlerule]
\titlespacing{\section}{0pt}{7pt}{3pt}

\setlist[itemize]{leftmargin=*, itemsep=0.5pt, topsep=1pt, parsep=0pt}
\setlist[enumerate]{leftmargin=*, itemsep=0.5pt, topsep=1pt, parsep=0pt}
\setlength{\parindent}{0pt}
\setlength{\parskip}{1pt}

\newcommand{\role}[2]{\textbf{#1} \hfill #2 \\}

\begin{document}
\pagestyle{empty}

%------------------------------------------------------------------ HEADER ----
\begin{center}
    {\LARGE\bfseries Musa Ibne Mannan}\\[5pt]
    Dallas, Texas 75252 \quad | \quad (737) 221-1884 \quad | \quad
    \href{mailto:musa314@utdallas.edu}{musa314@utdallas.edu}\\[3pt]
    \href{https://musamannan.com}{musamannan.com} \quad | \quad
    \href{https://www.linkedin.com/in/musa314/}{linkedin.com/in/musa314} \quad | \quad
    \href{https://orcid.org/0009-0005-5579-6899}{ORCID 0009-0005-5579-6899}
\end{center}
\vspace{-4pt}
\hrule height 1.5pt
\vspace{6pt}

%----------------------------------------------------------------- SUMMARY ----
\section{Summary}
Mechanical engineering Ph.D. candidate working on electrochemical biosensors, CVD synthesis of 2D materials, and
machine learning for process transfer. Built and commissioned a dual-zone CVD line from an empty bench, ran 139
MoS$_2$ growths, and trained a Random Forest on calibrated substrate temperature that reached 89.3\% accuracy and
carried a working recipe into a different reactor on the first attempt. Co-developed a screen-printed carbon
biosensor detecting tRNA-derived fragments at a 0.4 nM limit of detection. Six journal papers and proceedings, two
as first author. Seeking device fabrication, process development, and materials machine learning roles.

%--------------------------------------------------------------- EDUCATION ----
\section{Education}
\role{Doctor of Philosophy (Ph.D.) in Mechanical Engineering}{Aug 2023 -- Present}
Erik Jonsson School of Engineering and Computer Science, The University of Texas at Dallas, Richardson, TX \\
Concentration: Manufacturing and Design Innovations \quad GPA: 3.62

\role{Master of Science in Mechanical and Manufacturing Engineering (ABET accredited)}{Aug 2021 -- May 2023}
Ingram School of Engineering, Texas State University, San Marcos, TX \quad GPA: 3.75 \\
Thesis advisor: Dr. Namwon Kim

%------------------------------------------------------------------ SKILLS ----
\section{Technical Skills}
\textbf{Failure Analysis and Testing:} Root cause analysis, fracture surface examination (SEM, EDAX), mechanical
testing (tensile, flexural, Izod impact; MTS 810), thermal analysis (TGA, DSC, TMA), design review, field error
resolution. \\
\textbf{Simulation and CAD:} Abaqus finite element analysis (transient heat transfer, thermal stress, plasticity from true stress-strain
data, mesh convergence, 2D plane stress and 3D solid elements), SolidWorks, AutoCAD, CATIA, Mastercam, basic CFD. \\
\textbf{Programming and Data Analysis:} Python (data analysis pipelines, Random Forest, XGBoost), MATLAB
(numerical methods, custom FEM implementation, Gauss quadrature, multi-objective genetic algorithms). \\
\textbf{Experimental and Fabrication:} High-vacuum chemical vapor deposition (CVD) design and commissioning, pulsed laser deposition, Raman,
XPS, XRD, SEM/EDX, ellipsometry, Hall measurement, screen-printed electrode functionalization (oxygen plasma,
EDC/NHS), cyclic and differential pulse voltammetry, FDM/SLA/Markforged printing, twin-screw extrusion, CNC
(HAAS VF-3), laser cutting, waterjet, TIG and arc welding, mill and lathe. \\
\textbf{Professional:} Technical report writing, publication-quality figures and presentations, EHS and OSHA
compliance, procurement and inventory management, Adobe Creative Suite.

%-------------------------------------------------------------- EXPERIENCE ----
\section{Professional Experience}

\role{Research Assistant and Lab In-Charge}{Jun 2024 -- Present}
Leem Research Group, The University of Texas at Dallas, Richardson, TX
\begin{itemize}
    \item Independently designed, assembled from the ground up, troubleshot, and operated high-vacuum CVD furnaces, a hydrogen generator, and integrated Raman spectroscopy instrumentation, then optimized the resulting process.
    \item Led complete laboratory build-out including procurement, equipment installation, safety compliance, and inventory management while running concurrent computational and experimental programs.
    \item Applied Python-based data analysis pipelines and ensemble machine learning (Random Forest, XGBoost) to optimize multi-parameter mechanical system performance and materials synthesis processes.
\end{itemize}

\role{Graduate Teaching Associate (Instructor of Record)}{Aug 2025 -- Present}
The University of Texas at Dallas, Richardson, TX
\begin{itemize}
    \item Instructor of record for ECS 1100 within a 300-student course, delivering curriculum, mentoring, and assessment for an assigned section; lecture on engineering fundamentals including mechanics and design. Cumulative teaching load across UT Dallas and Texas State exceeds 800 students.
\end{itemize}

\role{Assistant and Site Engineer}{Sep 2018 -- Jun 2019}
Washillah Engineering Services (Partex Petro Ltd. and CUFL projects), Bangladesh
\begin{itemize}
    \item Led on-site teams of 28 to 35 engineers and contractors for large-scale mechanical and piping projects (condensate refinery and boiler re-tubing); enforced EHS and OSHA standards and managed field operations under tight deadlines.
    \item Performed root cause analysis and resolved over 120 design and implementation errors in mechanical systems and structures, saving approximately 30,000 man-hours and \$1M while maintaining quality and schedule.
    \item Coordinated cross-functional teams, client interfaces, material logistics, and quality control in high-stakes industrial field environments.
\end{itemize}

\role{Graduate Instructional Assistant}{Aug 2023 -- May 2024}
The University of Texas at Dallas, Richardson, TX
\begin{itemize}
    \item Graded assignments and exams and tutored 136 students in Kinematics and Dynamics (MECH 3350) and approximately 60 in Mechanics of Materials (MECH 2320).
    \item Delivered substitute lectures and assisted students with computational modeling, FEA, and CAD troubleshooting.
\end{itemize}

\role{Graduate Instructional Assistant}{Sep 2021 -- May 2023}
Texas State University, San Marcos, TX
\begin{itemize}
    \item Mentored three senior design teams (15 students) to completion on industry-sponsored mechanical engineering capstone projects; provided hands-on guidance in CAD, FEA, prototyping, and mechanical system design.
\end{itemize}

%---------------------------------------------------------------- PROJECTS ----
\section{Selected Projects}
\begin{itemize}
    \item \textbf{Machine Gun Barrel Thermal and Stress Failure Analysis (Abaqus FEA):} Calculated firing pressure analytically at 130.5 MPa, then ran transient heat transfer and von Mises stress analysis under sustained high-rate firing; evaluated AISI 4140, 316L, and ASTM A656 Grade 7 for thermal gradients and failure risk.
    \item \textbf{Cantilever Beam FEA and Mesh Convergence Study (Abaqus, MATLAB):} CPS4, CPS8, and C3D8 elements across three mesh densities; 3D deflection error fell from 84\% to 22\%, 2D mid-span stress error from 94\% to 8.6\%. Verified against custom MATLAB code using 2x2 Gauss quadrature.
    \item \textbf{Biocomposite Characterization and Failure Diagnosis:} Designed, extruded, and 3D-printed silane-treated natural fiber composites; used SEM/EDAX, tensile and flexural testing, and thermal analysis to diagnose interfacial adhesion issues, fiber-size effects, and resulting mechanical failure modes.
    \item \textbf{rGO Thin-Film Processing and Characterization:} Independently designed and executed pulsed laser deposition and annealing experiments; performed full characterization (XPS, Raman, XRD, ellipsometry, Hall measurement, low-temperature transport).
\end{itemize}

%------------------------------------------------------------ PUBLICATIONS ----
\section{Publications}
\begin{enumerate}
    \item \textbf{Musa Ibne Mannan}, Varin Kansal, Premesh N. Mistry, et al. ``Closing the Reproducibility Gap in MoS$_2$ CVD: Zero-Shot Cross-Reactor Transfer of Machine-Learned Growth Conditions via Thermal Calibration.'' \textit{Advanced Materials} (Wiley). Under review.
    \item \textbf{Musa Ibne Mannan}, Premesh Mistry, Varin Kansal, et al. ``AI-Enabled CVD Synthesis, Quality Control, and Autonomous Manufacturing of 2D Materials.'' \textit{npj Advanced Manufacturing} (Springer Nature), 2026. DOI: 10.1038/s44334-026-00109-5.
    \item Wasi Shadman, \textbf{Musa Mannan}, Tanzila Kamal Choity, et al. ``Electrochemical Detection of tRNA-Derived Fragment (tRF) Biomarkers Using a Multi-Functionalized Carbon Electrode.'' \textit{Surfaces and Interfaces} 83, 108446, 2026. DOI: 10.1016/j.surfin.2026.108446.
    \item Mohammad Shahed H.K. Tushar, \textbf{Musa Ibne Mannan}, Afra Azmain. ``Exhaust Gas After Treatment using Air Preheating and Selective Catalytic Reduction by Urea to Reduce NOx in Diesel Engine.'' \textit{Heliyon} 11(1), e42399, 2025. DOI: 10.1016/j.heliyon.2025.e42399.
    \item Jitendra S. Tate, Subash Pant, \textbf{Musa I. Mannan}, et al. ``Thermal and Structural Properties of Surface Treated Ragweed Fiber Reinforced Biodegradable Polylactic Acid Biocomposites.'' \textit{International Journal of Composite and Constituent Materials} 8(2), 2022.
    \item S. Panta, \textbf{M.I. Mannan}, et al. ``Comparison of Mechanical and Morphological Properties of Ameo and Glymo Silane Treated Ragweed-PLA Bio-Composites.'' \textit{CAMX 2022 Proceedings}, Anaheim, CA, 2022.
\end{enumerate}

%----------------------------------------------------- ADDITIONAL EXPERIENCE --
\section{Additional Experience}
\role{Medical Scribe and ER Workflow Innovation, Augmedix}{Jul 2019 -- Jul 2021}
Supported a California emergency room through COVID, delivering real-time documentation for 4,000+ patients. Designed and piloted the world's first remote ER scribing program, building the training curriculum from scratch and deploying the workflow across multi-site teams. Featured in \textit{Business Insider}.

\role{Novelist and Screenwriter (Bangla)}{Feb 2016 -- Present}
Author of 26 published books across Bangladesh and India, written as Kishor Pasha Imon alongside full-time engineering and research work.

\end{document}
